% Options for packages loaded elsewhere
\PassOptionsToPackage{unicode}{hyperref}
\PassOptionsToPackage{hyphens}{url}
\documentclass[
]{article}
\usepackage{xcolor}
\usepackage{amsmath,amssymb}
\setcounter{secnumdepth}{-\maxdimen} % remove section numbering
\usepackage{iftex}
\ifPDFTeX
  \usepackage[T1]{fontenc}
  \usepackage[utf8]{inputenc}
  \usepackage{textcomp} % provide euro and other symbols
\else % if luatex or xetex
  \usepackage{unicode-math} % this also loads fontspec
  \defaultfontfeatures{Scale=MatchLowercase}
  \defaultfontfeatures[\rmfamily]{Ligatures=TeX,Scale=1}
\fi
\usepackage{lmodern}
\ifPDFTeX\else
  % xetex/luatex font selection
\fi
% Use upquote if available, for straight quotes in verbatim environments
\IfFileExists{upquote.sty}{\usepackage{upquote}}{}
\IfFileExists{microtype.sty}{% use microtype if available
  \usepackage[]{microtype}
  \UseMicrotypeSet[protrusion]{basicmath} % disable protrusion for tt fonts
}{}
\makeatletter
\@ifundefined{KOMAClassName}{% if non-KOMA class
  \IfFileExists{parskip.sty}{%
    \usepackage{parskip}
  }{% else
    \setlength{\parindent}{0pt}
    \setlength{\parskip}{6pt plus 2pt minus 1pt}}
}{% if KOMA class
  \KOMAoptions{parskip=half}}
\makeatother
\usepackage{graphicx}
\makeatletter
\newsavebox\pandoc@box
\newcommand*\pandocbounded[1]{% scales image to fit in text height/width
  \sbox\pandoc@box{#1}%
  \Gscale@div\@tempa{\textheight}{\dimexpr\ht\pandoc@box+\dp\pandoc@box\relax}%
  \Gscale@div\@tempb{\linewidth}{\wd\pandoc@box}%
  \ifdim\@tempb\p@<\@tempa\p@\let\@tempa\@tempb\fi% select the smaller of both
  \ifdim\@tempa\p@<\p@\scalebox{\@tempa}{\usebox\pandoc@box}%
  \else\usebox{\pandoc@box}%
  \fi%
}
% Set default figure placement to htbp
\def\fps@figure{htbp}
\makeatother
\setlength{\emergencystretch}{3em} % prevent overfull lines
\providecommand{\tightlist}{%
  \setlength{\itemsep}{0pt}\setlength{\parskip}{0pt}}
\usepackage{bookmark}
\IfFileExists{xurl.sty}{\usepackage{xurl}}{} % add URL line breaks if available
\urlstyle{same}
\hypersetup{
  pdftitle={01-04\_Tag\_4\_Finalisierung\_Uebergabe},
  hidelinks,
  pdfcreator={LaTeX via pandoc}}

\title{01-04\_Tag\_4\_Finalisierung\_Uebergabe}
\author{}
\date{}

\begin{document}
\maketitle

\subsection{Tag 4: Finalisierung \&
Übergabe-Vorbereitung}\label{tag-4-finalisierung-uxfcbergabe-vorbereitung}

\textbf{Datum:} 19.02.2026\\
\textbf{Status:} ✅ Dokumentation abgeschlossen \textbar{} ✅
Visualisierung abgeschlossen

\subsubsection{Tagesziele (PM-4)}\label{tagesziele-pm-4}

Lösungsansätze finalisieren und bewerten (Entscheidung für Lösung A
gefallen)

Arbeitsplanung für die kommenden Wochen (Themenfeld 2 strukturiert)

Übergabe-Dokumentation erstellen (Finaler Entwurf)

Packet Tracer Datei (KL) sichten \& detailliert analysieren

Organisationsstruktur visualisieren (Organigramm integriert)

\begin{center}\rule{0.5\linewidth}{0.5pt}\end{center}

\subsection{1. Arbeitsplanung Netzwerkadministration
(Update)}\label{arbeitsplanung-netzwerkadministration-update}

Das Netzwerk-Team erhält von uns folgendes \textbf{finales Paket} für
die Umsetzung in Themenfeld 2:

\begin{itemize}
\item
  \textbf{VLAN-IDs \& Namen:}\\
  10 -- Vertrieb\\
  20 -- IT\\
  30 -- HR\\
  40 -- VoIP\\
  99 -- Management
\item
  \textbf{IP-Adressbereiche:} 10.13.X.0/24 (X = VLAN-ID)
\item
  \textbf{DHCP-Strategie:} Zentrales Management über Admin-VM
  (10.8.13.2) via Relay-Agent auf allen SVIs
\item
  \textbf{Hardware-Set:}\\
  2× Catalyst 3650 (Core -- Stacking + HSRP-Redundanz geplant)\\
  3× Catalyst 2960 (Access)\\
  1× pfSense (Mainz -- VPN-Termination)\\
  \pandocbounded{\includegraphics[keepaspectratio,alt={Screenshot 2026-02-20 120006.png}]{C:/Users/User/Documents/SamuelDenkwerkstätte/Screenshot 2026-02-20 120006.png}}
\end{itemize}

\begin{center}\rule{0.5\linewidth}{0.5pt}\end{center}

\subsection{2. Finale Übergabe-Dokumentation
(Struktur)}\label{finale-uxfcbergabe-dokumentation-struktur}

\subsubsection{Projektüberblick}\label{projektuxfcberblick}

Modernisierung der BetaTrade-Infrastruktur. Kernstück ist der Ausbau
Kaiserslauterns zum Regional-Hub mit Hochverfügbarkeits-Ansatz (HSRP,
redundante Trunks).

\subsubsection{Priorisierte Anforderungen \&
Erkenntnisse}\label{priorisierte-anforderungen-erkenntnisse}

\begin{enumerate}
\item
  \textbf{Sicherheit}\\
  Isolation der Abteilungen durch VLANs. Port-Security (802.1X) auf
  Access-Switches für IT-Bereich zwingend erforderlich.
\item
  \textbf{Konnektivität}\\
  S2S-VPN (Site-to-Site) zwischen Mainz und KL. Zertifikatsbasierte
  Authentifizierung statt Passwort.
\item
  \textbf{Dienste}\\
  DHCP-Relay muss auf den Core-Switches für alle Client-VLANs aktiv sein
  (\texttt{ip\ helper-address\ 10.8.13.2}).
\end{enumerate}

→ Detaillierte Übergabe: Uebergabe\_Netzwerk\_Team

\subsubsection{Ausgearbeitete Lösung (A --
Security-Fokus)}\label{ausgearbeitete-luxf6sung-a-security-fokus}

Wir implementieren \textbf{Lösung A}. Diese beinhaltet:

\begin{itemize}
\tightlist
\item
  VLAN-Segmentierung + Inter-VLAN-Routing nur über Core
\item
  802.1X Port-Security (IT-VLAN)
\item
  pfSense OpenVPN mit Zertifikaten
\item
  HSRP-Redundanz zwischen beiden Core-Switches
\item
  QoS für VoIP (VLAN 40 priorisieren)
\item
  Vorbereitung für IDS/IPS in Phase 3
\end{itemize}

\begin{center}\rule{0.5\linewidth}{0.5pt}\end{center}

\subsection{3. Detaillierte Analyse Packet Tracer
(Kaiserslautern\_initial.pkt)}\label{detaillierte-analyse-packet-tracer-kaiserslautern_initial.pkt}

\textbf{Inventar-Check:}\\
Die physische Topologie entspricht dem hierarchischen Design (Core →
Distribution → Access).\\
Neu: Es sind bereits IP-Telefone und PCs an den Access-Switches
angeschlossen, jedoch ohne Link-Licht (Ports teilweise noch im
\texttt{shutdown}).

\textbf{Konfigurations-Lücken (zu erledigen in TF-2):}

\begin{itemize}
\tightlist
\item
  \textbf{VTP/VLANs:} Noch keine VLAN-Datenbanken vorhanden. Empfehlung:
  VTP-Modus \texttt{transparent} nutzen.
\item
  \textbf{STP (Spanning Tree):} Aufgrund redundanter Anbindung →
  \texttt{Rapid-PVST} konfigurieren, um Loops zu verhindern.
\item
  \textbf{Layer 3:} Core-Switches laufen aktuell nur Layer 2. Erster
  Schritt: \texttt{ip\ routing} aktivieren.
\item
  \textbf{VoIP:} Telefone benötigen Voice-VLAN (40) + PoE.
  Leistungsbilanz der Switches prüfen (PoE-Budget).
\end{itemize}

\textbf{Offene Fragen aus PT-Analyse:}

\begin{itemize}
\tightlist
\item
  Management-VLAN 99: Separates Subnetz außerhalb 10.13.X.X sinnvoll?
\item
  Voice-VLAN: Inline-Power + Data-VLAN durchschleifen (PC hinter
  Telefon) korrekt unterstützt?
\end{itemize}

\begin{center}\rule{0.5\linewidth}{0.5pt}\end{center}

\subsection{4. Organisationsstruktur
(Organigramm-Erkenntnisse)}\label{organisationsstruktur-organigramm-erkenntnisse}

Das Organigramm wurde erstellt und zeigt die strikte Trennung:

\begin{itemize}
\tightlist
\item
  \textbf{Zentrale Mainz:} Management \& Kern-Server (AD, DB,
  pfSense).\\
\item
  \textbf{Regional-Hub KL:} Operatives Geschäft (Vertrieb, HR) mit
  lokaler IT-Präsenz.\\
\item
  \textbf{Infrastruktur-Diagramm:} Ein zweites Diagramm wurde erstellt,
  das die logische Zuordnung der Server in Mainz zu den neuen VLANs in
  KL
  visualisiert.\pandocbounded{\includegraphics[keepaspectratio,alt={Screenshot 2026-02-20 120542.png}]{C:/Users/User/Documents/SamuelDenkwerkstätte/Screenshot 2026-02-20 120542.png}}
\end{itemize}

\begin{center}\rule{0.5\linewidth}{0.5pt}\end{center}

\subsection{5. Offene Fragen \& Risikomanagement (Update Tag
4)}\label{offene-fragen-risikomanagement-update-tag-4}

\begin{itemize}
\tightlist
\item
  \textbf{Risiko (VoIP VLAN 40):} PCs hinter Telefonen → korrekte
  Voice-VLAN-Zuweisung prüfen (Packet Tracer Test)
\item
  \textbf{Risiko (Hardware/IOS):} Catalyst 3650 benötigen passendes IOS
  für Inter-VLAN-Routing + HSRP
\item
  \textbf{Frage an Technik/Kunde:} Management-Zugriff (VLAN 99) →
  separates Subnetz (z. B. 10.99.99.0/24) oder innerhalb 10.13.99.0/24
  belassen?
\end{itemize}

→ Verlinkte Dokumente:\\
Uebergabe\_Netzwerk\_Team \textbar{} Risiko-Register BetaTrade
\textbar{} Phase 1 Kaiserslautern \textbar{} BetaTrade Modernisierung

\textbf{Letzte Aktualisierung:} 19.02.2026 -- Projektassistenz\\
\textbf{Nächster Schritt:} Übergabe-Meeting Themenfeld 2 (KW 8)

\section{Tag 4: Visualisierung \& Handover-Vorbereitung
(Ergänzung)}\label{tag-4-visualisierung-handover-vorbereitung-erguxe4nzung}

\begin{quote}
\textbf{INFO:} Fokus\\
Grafische Aufbereitung der Netzwerkstruktur und formale Vorbereitung der
Übergabe an das IT-Implementierungsteam.
\end{quote}

\subsection{1. Diagramm-Integration und
Workflow}\label{diagramm-integration-und-workflow}

\begin{itemize}
\tightlist
\item
  \textbf{Prozess:} Dokumentation, wie logische Netzwerktopologien
  (Packet Tracer Exporte, D2-Skripte) versioniert und als Bild-Assets
  sauber in den Markdown-Text eingebunden werden, ohne den Lesefluss zu
  stören.
\end{itemize}

\subsection{2. Roadmap-Formatierung}\label{roadmap-formatierung}

\begin{itemize}
\tightlist
\item
  Aufbereitung der mehrstufigen Modernisierungs-Roadmap für die Zentrale
  Mainz in einem präsentationstauglichen Format für das
  Management-Meeting.
\end{itemize}

\subsection{3. Übergabe-Checkliste
(Handover)}\label{uxfcbergabe-checkliste-handover}

\begin{quote}
\textbf{CHECK:} Handover-Status\\
Formale Checkliste erstellt, um sicherzustellen, dass das Netzwerk-Team
starten kann:

VLAN-IDs und Namenskonventionen definiert.

IP-Adressbereiche (Subnetze) zugewiesen.

DHCP-Relay-Strategie dokumentiert.
\end{quote}

\end{document}
